\chapter{Cronograma}

\textcolor{red}{Cronograma provisório}

\begin{table}[H]
	\centering
	\footnotesize
	\caption{Cronograma previsto para o desenvolvimento da CNN1D e síntese com FINN (Fev–Out/2026).}
	\renewcommand{\arraystretch}{1.3}
	\begin{tabular}{p{6cm} c c c c c c c c c}
		\toprule
		\textbf{Atividade} &
		\textbf{Fev} &
		\textbf{Mar} &
		\textbf{Abr} &
		\textbf{Mai} &
		\textbf{Jun} &
		\textbf{Jul} &
		\textbf{Ago} &
		\textbf{Set} &
		\textbf{Out} \\
		\midrule
		Preparação e limpeza da base de dados & X & X &   &   &   &   &   &   &   \\
		Segmentação e definição das janelas (pré-processamento) & X & X &   &   &   &   &   &   &   \\
		\midrule
		Definição da arquitetura CNN1D &   & X & X &   &   &   &   &   &   \\
		Treinamento inicial e ajuste de hiperparâmetros &   &   & X & X &   &   &   &   &   \\
		Validação da CNN1D (métricas clássicas) &   &   &   & X & X &   &   &   &   \\
		\midrule
		Quantização com Brevitas (QAT) &   &   &   & X & X &   &   &   &   \\
		Exportação para QONNX e depuração &   &   &   &   & X & X &   &   &   \\
		\midrule
		Geração do pipeline \textit{dataflow} no FINN &   &   &   &   &   & X & X &   &   \\
		Simulação de latência, recursos e throughput &   &   &   &   &   & X & X &   &   \\
		Ajuste fino e iteração entre FINN e modelo &   &   &   &   &   &   & X & X &   \\
		\midrule
		Avaliação final dos resultados (software + FINN) &   &   &   &   &   &   &   & X & X \\
		Redação dos resultados e consolidação do capítulo &   &   &   &   &   &   &   &   & X \\
		\bottomrule
	\end{tabular}
\end{table}
